% !TEX root = lectures.tex
\newcommand{\Cut}{\textrm{Cut}}
\section{Lecture 1}
In this class, we will explore the following themes.
\begin{enumerate}
    \item Randomness - When the worst-case model means there are not good algorithms,
    randomized error models can give you better guarantees, in line with reality. Randomized
    algorithms also can be better (on average) than deterministic ones on worst-case inputs.
    \item Optimization - Continuous optimization
    can help design great algorithms. Examples include linear programming, semidefinite programming,
    gradient descent, and stochastic gradient descent.
    \item Algebraic Techniques - Singular Value Decomposition, low-rank approximation
    and spectral methods use linear algebra to build useful algorithms. Furthermore, we
    can try to solve systems of polynomial equations.
\end{enumerate}
Even more important is the interplay between these three themes.

\subsection{Randomized Min Cuts}
Karger's 1993 paper gave a simple algorithm for finding minimum cuts in graphs.
Define a graph $G$ with vertices $V$ and edges $E$. Here $ n = |V|$ and $m = |E|$.
\begin{definition}
    Given a graph $G = (V, E)$ and $S \subseteq V$,
    define the cut $\Cut(S) = E(S, \bar{S})$.
    Our goal is to find $S \neq \emptyset, V$ that minimizes $|\Cut(S)|$.
\end{definition}

\begin{definition}[Contraction]
    The contraction of an edge $e \in E$ produces a new graph $G' = (V', E)$
    where $V' = (V \setminus \{u, v\}) \cup S_{u, v}$ (combine $u$ and $v$ into a ``super-vertex'') and
    \[ E' = \{ \{S_{u, v}, w\} \mid \{u, w\} \text{ or } \{v, w\} \in E \} \cup \{ e \in E \mid u \notin e, v \notin e \} \]
    where we take the edges that we add as a multigraph (multiple edges between the same pair of vertices).
\end{definition}

Notice that every time a contraction is performed,
the number of vertices decreases by 1 and the number of edges decreases
(by maybe more than one, due to multiedges). Here is the contraction algorithm.

\begin{algothm}[Karger's Min Cut]
    \begin{itemize}
        \item $H = G$, $H = (V', E')$.
        \item While $H$ has more than $2$ vertices:
        \item \,\,\,\,\, Pick a random edge on $E'$, $\{u, v\}$.
        \item \,\,\,\,\, Contract $H$ on $\{u, v\}$.
        \item Consider the number of edges between the two vertices
        left as a minimum cut candidate.
    \end{itemize}
\end{algothm}

\begin{theorem}
    Fix a min-cut $(X, \bar{X})$ in $G$. We claim the algorithm above returns $|\Cut(X)|$ with probability at least
    $\frac{2}{n(n-1)}$.
\end{theorem}

Thus, it suffices to run the original algorithm independently $t$ times and take the smallest cut we see.
The probability that we make a mistake all $t$ times is
\[ \P{\text{failure chance}} \le \qty(1  - \frac{2}{n(n - 1)})^{t} \le e^{-2\frac{t}{n(n-1)}} \]
if we take $t = O(n^2 \log n)$, then the failure chance is at most $\frac{1}{n^{O(1)}}$. Furthermore,
the contractions each take $O(n)$ time and there are $O(n)$ iterations. Thus, the runtime of the whole
algorithm is poly-time, taking at most $O(n^4 \log n)$ time.
We will soon find a speedup, but first let us prove
the theorem.

\begin{proof}
    Note that any cut in the contracted graph $\Cut(S)$ in $H$ can be mapped one-to-one to
    some cut in $G$. Further, every cut in $G$ that doesn't split a contracted pair of vertices
    has a corresponding cut in $H$. Furthermore, a cut survives a step of Karger's if you do not contract
    one of its edges. The probability that you miss the cut $(X, \bar{X})$
    is $1 - \frac{|\Cut(X)|}{|E'|}$. Furthermore, suppose the minimum cut has $c$ edges,
    then $G$ has at least $\frac{cn}{2}$ edges in total. This is because the degree of each 
    vertex must be at least
    $c$. Thus, at every step, the survival probability is at least $1 - \frac{c}{\frac{c|V'|}{2}}$.
    Thus, overall
    \[\P{\text{survival of cut}}  \ge \prod_{i = 3}^{n} \qty(1 - \frac{2}{i}) =  \prod_{i = 3}^{n} \qty(\frac{i - 2}{i}) = \frac{(n - 2)!}{n!/2!} = \frac{2}{n(n-1)} \]
\end{proof}

\subsection{Improving Min Cut}
We saw that when the number of vertices is large, it's unlikely that we contract the min cut.
So, we should not repeat the beginning too many times.
\begin{theorem}
    The chance that the minimum cut survives $n - \frac{n}{\sqrt{2}}$ steps of contraction
    is $\ge \frac{1}{2}$. 
\end{theorem}
Now, we can simply do two independent contractions until there are $\frac{n}{\sqrt{2}}$ vertices in each.
Then we can recursively calculate minimum cuts in both and take the better one. Then
we can use a recurrence to calculate the running time.
\[ T(n) = O(n^2) + 2T\qty(\frac{n}{\sqrt{2}}) \implies T(n) = O(n^2 \log n) \]
Finally, we should prove the following theorem
\begin{theorem}
    The algorithm finds a minimum cut with probability at least $\frac{\Delta}{\log_2 n}$.

    \begin{proof}
        Call $P(n)$ the chance the algorithm works on a graph of size $n$. Then $P(n) \ge 1 - (1 - \frac{1}{2} P(n/\sqrt{2}))^2$.
        By induction, we can prove the claim.
    \end{proof}
\end{theorem}